\section{Method}\label{sec:method}

The design below was committed to the repository on 29 October 2026, before
any loan-level file had been downloaded (commit \texttt{07c8ffd}). The data log,
the deviations from it, the corrected bugs and the analyses added after the
results are listed in \cref{app:deviations}. Each choice is
stated with the reason for it.

\subsection{Reporting dates, portfolio and training data}

A reporting date $T$ is a calendar year-end, December 2005 to December 2025
(\num{const.n_dates} dates). The \emph{portfolio} at $T$ is every loan active
at December $T$: balance above zero, no zero-balance event, and never yet 180
or more days delinquent. Loans 30 to 179 days past due at $T$ are in the
portfolio, but current delinquency is not a model input (the design has
none), so the models treat them like current loans with the same inputs;
\cref{sec:limits} reports how much of the book they are. Allowances are computed on the \num{const.n_port}
portfolio loans with the smallest value of a fixed loan-level random draw,
and expressed as a rate, expected loss over balance on that sample (a ratio
estimator). The dollar figures multiply the rate by the balance of the whole
active book. \Cref{sec:sampling} reports the sampling error this adds.

A model at $T$ may use only what a lender knew at $T$, in the data as
published in 2026: Freddie Mac's Release 47 and current FRED vintages, which
may revise what was reported at the time (\cref{sec:limits}). Its training data are
the \emph{loan-years} whose outcome year ends on or before $T$: a loan active
at December $y \le T-1$ contributes one row, with inputs measured at December
$y$ and the outcome observed in year $y+1$. Censored loan-years (removal for a
defect or a reperforming sale, without a default first) are dropped. Loans,
never rows, are split by a keyed hash of the loan identifier: 20\% are held
out and used only to measure accuracy, 10\% are used for early stopping and
70\% for fitting. Each loan is kept with probability
$\min\{1, 6.25\text{M}/n_T\}$, where $n_T$ is the number of loan-years
available, so that fitting and early-stopping rows together are capped near
the design's five million.

\subsection{Outcomes, competing risks and inputs}

In each loan-year a loan defaults, prepays or survives. \emph{Default} is the
first month 180 or more days delinquent, or a credit-event zero balance
(third-party sale, short sale or charge-off, REO disposition, whole-loan
sale; codes 02, 03, 09, 15), whichever comes first. \emph{Prepayment} is a
voluntary payoff (code 01). Two binary hazard models describe the year:
\begin{equation}
h_d = \Pr(\text{default}\mid\text{active}),\qquad
h_p = \Pr(\text{prepay}\mid\text{active, no default}),
\end{equation}
so that $\Pr(\text{default})=h_d$, $\Pr(\text{prepay})=(1-h_d)h_p$ and
$\Pr(\text{survive})=(1-h_d)(1-h_p)$. Any three-outcome distribution factors
this way, and the multinomial log loss is exactly the sum of the two binary log
losses, so nothing is lost relative to a multinomial model. A single
multiclass gradient-boosting model was rejected because LightGBM applies a
monotone constraint in the same direction to every class's score, which would
force prepayment to rise with leverage as default does.

The \num{const.n_features} inputs are those in \cref{tab:features}: credit
score, original LTV and combined LTV, debt-to-income ratio, original rate and
its spread over the Freddie Mac PMMS 30-year rate in the origination month,
loan purpose, occupancy, property type, units, first-time-buyer flag, number
of borrowers, origination channel, mortgage-insurance percentage, original
term, loan age, \emph{current} LTV (current balance over the origination value
marked to the FHFA state house-price index), the state's house-price growth
over the prior year, and its unemployment rate and twelve-month change. Seller
and servicer names, ZIP3 and MSA are not used: they identify institutions or
come close to identifying places.

\subsection{The allowance}

For each portfolio loan the projection runs one row per remaining contractual
year $k=1,\dots,K$. The balance at the start of year $k$, $B_k$, is the
interest-bearing balance amortised at the loan's current rate over its
remaining term, plus any non-interest-bearing (deferred) balance, which is due
at maturity. Loan age advances by twelve months a year, current LTV follows
$B_k$ and the house-price path, and the origination inputs do not change.
With $S_1=1$,
\begin{equation}\label{eq:el}
m_k = S_k\,h_{d,k},\qquad S_{k+1}=S_k\,(1-h_{d,k})(1-h_{p,k}),\qquad
\mathrm{EL} = \sum_{k=1}^{K} B_k\, m_k\, \mathrm{LGD}_k ,
\end{equation}
where a final part-year of $q$ months uses $1-(1-h)^{q/12}$. The allowance is
undiscounted, the probability-of-default/loss-given-default method, which ASC
326 permits without discounting outside a discounted-cash-flow approach. The
lifetime PD of a loan is $\sum_k m_k$.

\paragraph{Macro path.} During the two-year reasonable and supportable period
the state's house-price growth and unemployment rate stay at their December $T$
values. Afterwards they revert to their long-run state means, computed over
all years up to $T$, in a straight line over one year (primary). The two
other policies of H2 are immediate reversion and straight-line reversion over
three years; ASC 326-20-30-9 names immediate and straight-line reversion, and
the durations are this study's choice. A loan-year beginning at year-end $p$
reads the path's value at $p$, as in training. Because the path is annual, the
first reverted value is at $T+3$ under both immediate and one-year
straight-line reversion, so those two policies give the same allowance and the
policy range of H2 is in effect the range between one- and three-year
reversion. The length of the reasonable and supportable period, also the
lender's choice, is held at two years throughout. A secondary
analysis replaces the forecast with the realized path through 2025.

\paragraph{Loss given default.} One severity model per date, held fixed across
all default models: total actual loss over total balance at default (a
balance-weighted severity), among defaults resolved by December $T$, by current-LTV band at the start of the
default year ($<$60, 60--80, 80--90, 90--100, 100--120, $\ge$120) and by
judicial-foreclosure state \citep{gerardi2013borrower}. A default
later paid off voluntarily has a loss of zero. A cell with fewer than thirty
resolutions uses its band's mean; a band with fewer than thirty uses the mean
of all resolutions. Actual loss is Freddie Mac's own calculation
\citep{FreddieGuide2026}.

\subsection{Models}

The default model is one of \num{const.n_models}, all trained on the same rows:
\begin{itemize}[nosep]
\item 32 LightGBM binary models \citep{ke2017lightgbm}: 15 or 63 leaves $\times$ four
row/column subsampling settings (80/100\% $\times$ 80/100\%), with five seeds
for each subsampled setting and one for the deterministic one;
\item the same 32 with monotone constraints: the default hazard may not fall as
current LTV or the debt-to-income ratio rises, or as the credit score falls;
\item one logistic regression on cubic splines of the numeric inputs (five
quantile knots) with one-hot categories.
\end{itemize}
LightGBM uses a learning rate of 0.1, at least 100 rows per leaf and early
stopping after 50 rounds without improvement on the early-stopping loans. The
prepayment model is fixed at each date: the most accurate on held-out
prepayment log loss among the eight unconstrained LightGBM settings (seed 1)
and the spline logit (Deviation~1).

\subsection{The set of equally accurate models}\label{sec:set}

Accuracy is held-out log loss of the default model, a proper scoring rule that
rewards calibration, which an allowance needs. Its sampling error is a
loan-clustered bootstrap: \num{const.n_boot} resamples of the held-out loans
with replacement. A model is in the set if its held-out log loss is within one
bootstrap standard error of the best. This is a tolerance,
not a test that the models are equal: a paired comparison on the same held-out
loans separates many of them from the best, and \cref{sec:sensitivity}
reports the much smaller sets such a test admits. The tolerance asks instead
which models a lender could not reject on accuracy measured as the lender
would usually measure it, one model at a time. Sensitivity analyses use half and two standard errors, and a set defined by
held-out AUC instead.

\subsection{Predictions}\label{sec:predictions}

Six predictions were recorded with the design, each with the result that
would refute it. They are reproduced here as written.

\begin{enumerate}[label=\textbf{H\arabic*.}, leftmargin=*, itemsep=2pt]
\item \emph{Model choice exceeds training noise.} At a majority of dates, the
range (max $-$ min) of the allowance across the set exceeds the 95\% interval
width of the best configuration's allowance across 20 refits on loan-level
bootstrap resamples of its training data.
\item \emph{Model choice rivals permitted policy discretion.} At a majority of
dates, the allowance range across the set is at least the range across the
three reversion policies ASC 326-20-30-9 names (immediate, straight-line over
one year, straight-line over three years), holding the best model fixed.
\item \emph{The spread widens into stress.} The allowance range as a share of
the best model's allowance is larger at year-ends 2006 and 2007 than at the
median of all other dates.
\item \emph{Constraints narrow the spread.} Within each family's own set (same
tolerance, relative to that family's best), the relative range is smaller for
monotone-constrained models than for unconstrained ones at a majority of dates.
\item \emph{The range is not forecast uncertainty.} For dates whose portfolio
has at least eight years of realized outcomes, realized lifetime credit losses
lie inside the range of allowances across the set at fewer than half of those
dates.
\item \emph{Staging flips.} Under the IFRS 9 analogue (lifetime PD at $T$ more
than twice the lifetime PD at origination as the significant-increase test),
the share of portfolio loans whose stage differs between at least two models
in the set exceeds the share whose stage differs across the 20 bootstrap
refits of H1, at a majority of dates.
\end{enumerate}

\noindent For H6 the lifetime PD at origination is computed by the same model
from the loan's origination inputs (age zero, original balance, rate and term,
current LTV equal to original LTV) and the macro values at the year-end before
its first payment, held for two years and then reverted as in the primary
path. A loan whose stage is the same under every model counts as stable. This
is a simplified analogue of the IFRS 9 test: it compares lifetime PDs from one
model, not the lender's own origination estimate; it compares the PD over the
remaining life with the PD over the full term, without the adjustment for the
shorter remaining life that IFRS 9 asks for (B5.5.11), which biases seasoned
loans towards stage 1; its origination path reverts to long-run means
computed at $T$; and it has no 30-days-past-due backstop.
